\title{Group Sequence Policy Optimization}

\begin{abstract}
We present Group Sequence Policy Optimization (GSPO), a reinforcement learning method characterized by its stability, efficiency, and effectiveness in training large language models. Departing from earlier approaches that employ importance ratios at the token level, GSPO calculates importance ratios using the likelihood of entire sequences, incorporating sequence-level clipping, reward calculation, and optimization steps. Our results show that GSPO not only outperforms the GRPO algorithm in terms of training efficiency and effectiveness, but also significantly enhances the stability of Mixture-of-Experts (MoE) RL training and offers prospects for streamlining RL system architecture. These advantages of GSPO have played a key role in driving the outstanding progress seen in the most recent Qwen3 models.
\end{abstract}

\section{Introduction}
\label{sec:intro}

Reinforcement learning (RL) has become a central framework in advancing the scalability of language models \citep{o1,dpsk-r1,qwq32b,qwen3}. By applying RL at large scales, these models acquire enhanced abilities to address complex tasks—including advanced mathematics and programming challenges—by engaging in more extended and intricate reasoning.

Ensuring stability and robustness in training dynamics is essential for effectively scaling RL with increased computational resources. Nevertheless, leading RL algorithms such as GRPO \citep{grpo} struggle with pronounced instability during the training of very large language models, frequently leading to catastrophic failures and permanent model collapse \citep{qwen3, minimax-m1}. Such instability presents a significant obstacle to advancing the performance of language models through extended RL optimization.

This work demonstrates that GRPO's instability arises primarily from the incorrect use and breakdown of importance sampling weights in its algorithm. Such misuse leads to elevated training noise with high variance, which intensifies as the response length grows and is exacerbated by the clipping process, eventually leading to model failure.

In response to these fundamental challenges, we introduce \textbf{Group Sequence Policy Optimization (GSPO)}, a novel reinforcement learning algorithm tailored for training large language models. GSPO's main contribution is its rigorously defined importance ratio, derived from sequence likelihood \citep{click}, in line with the foundational concept of importance sampling. Furthermore, GSPO determines the normalized rewards as the advantage values among several responses to the same query, thereby facilitating coherence between rewards at the sequence level and the optimization process.

Our experimental results highlight the considerable advantages of GSPO compared to GRPO regarding stability during training, computational efficiency, and overall performance. Notably, GSPO intrinsically overcomes the stability issues that typically arise in RL training for large Mixture-of-Experts (MoE) architectures, thereby obviating the necessity for intricate stabilization techniques and suggesting the possibility of streamlining RL system design. These benefits of GSPO have played a pivotal role in achieving the substantial performance gains observed in the latest Qwen3 models. We propose GSPO as a dependable and scalable algorithmic basis to support ongoing progress in large-scale RL training for language models.

\section{Preliminaries}

\paragraph{Notation}
Throughout this work, we represent an autoregressive language model with parameters $\theta$ as a policy $\pi_\theta$. The variable $x$ is used to signify a query, while $\mathcal{D}$ refers to the set of all queries. For a response $y$ corresponding to a particular query $x$, the probability assigned to $y$ by the policy $\pi_\theta$ is expressed as $\pi_\theta (y | x)=\prod_{t=1}^{|y|} \pi_\theta (y_t | x, y_{<t} )$, where $|y|$ indicates the total number of tokens in $y$. To assess the quality of a query-response pair $(x, y)$, a verifier $r$ produces a score, giving a reward $r(x, y)$ that falls within the interval $[ 0, 1 ]$.

\paragraph{Proximal Policy Optimization (PPO)} PPO \citep{ppo} utilizes samples collected from the previous policy $\pi_{\theta_\text{old}}$ and enforces that subsequent policy updates remain close to this old policy by leveraging a clipping strategy. More precisely, the algorithm optimizes the policy via the following objective function (excluding the KL regularization term for succinctness, since it is outside the scope of this study):

\begin{align}
\mathcal{J}_\text{PPO}(\theta) = \mathbb{E}_{ x \sim \mathcal{D},\, y \sim \pi_{\theta_\text{old}}( \cdot | x) }
\left[ \frac{1}{|y|} \sum_{t=1}^{|y|} 
\min \left( w_{t}(\theta) \widehat{A}_{t},  \, \mathrm{clip} \left( w_{t}(\theta), 1 - {\varepsilon}, 1 + {\varepsilon}\right) \widehat{A}_{t} \right)
\right],
\end{align}

Here, the importance ratio for the token $y_t$ is given by $
w_{t}(\theta) = \frac{ \pi_{\theta} (y_{t} | x, y_{<t}) }{ \pi_{\theta_\text{old}} (y_{t} | x,y_{<t})} $, with the advantage $\widehat{A}_{t}$ of $y_t$ being approximated using a separate value model. The parameter $\varepsilon$ determines the clipping threshold applied to the importance ratios.

A principal difficulty encountered when using PPO in real-world settings stems from its significant dependence on the value model. In practice, this model typically matches the policy model in scale, resulting in substantial demands on memory and computational resources. Moreover, the success of the algorithm critically depends on the precision of its value estimates. Not only is constructing an accurate value model inherently non-trivial, but extending its robustness to support longer outputs and more intricate tasks amplifies the complexity of the problem.

\paragraph{Group Relative Policy Optimization (GRPO)} GRPO \citep{grpo} circumvents reliance on a value model by evaluating the relative advantage among responses belonging to the same group, all generated in response to an identical query. More precisely, GRPO aims to optimize the following objective:

\begin{align}
\label{equ:grpo}
\mathcal{J}_\text{GRPO}(\theta) = \mathbb{E}_{ x \sim \mathcal{D},\, \{y_i\}_{i=1}^G \sim \pi_{\theta_\text{old}}( \cdot | x) }
\left[ \frac{1}{G} \sum_{i=1}^{G} \frac{1}{|y_i|} \sum_{t=1}^{|y_i|} 
\min \left( w_{i,t}(\theta) \widehat{A}_{i,t},  \, \mathrm{clip} \left( w_{i,t}(\theta), 1 - {\varepsilon}, 1 + {\varepsilon}\right) \widehat{A}_{i,t} \right)
\right],
\end{align}

Here, $G$ denotes the total number of responses produced for every input $x$ (that is, the size of the group). The importance ratio $w_{i,t}(\theta)$ and the advantage $\widehat{A}_{i,t}$ associated with the token $y_{i,t}$ are given as follows:

\begin{align}
    w_{i,t}(\theta)=\frac{ \pi_{\theta} (y_{i,t} | x, y_{i,<t}) }{ \pi_{\theta_\text{old}} (y_{i,t} | x,y_{i,<t})},\quad\ 
    \widehat{A}_{i,t} = \widehat{A}_{i} = \frac{r(x, y_i) - \mathrm{mean} \left( \{ r(x, y_i) \}_{i=1}^G \right) }{ \mathrm{std} \left( \{ r(x, y_i) \}_{i=1}^G \right) },
\end{align}

In this case, each token within $y_i$ is assigned the identical advantage value $\widehat{A}_{i}$.

\section{Motivation}
\label{sec:motivation}

The increase in model parameters, the adoption of sparsity patterns (such as those found in Mixture-of-Experts architectures), and longer response outputs all require the use of large rollout batch sizes to fully leverage hardware during reinforcement learning. In pursuit of greater sample efficiency, it is common to divide these large rollout batches into several smaller mini-batches for performing gradient updates. However, this approach naturally gives rise to an off-policy learning scenario, since the generated responses $y$ originate from a previous policy $\pi_{\theta_\text{old}}$ instead of the current policy $\pi_\theta$ targeted during optimization. This dynamic underscores the importance of employing clipping strategies in algorithms like PPO and GRPO, which serve to limit the influence of samples that are substantially "off-policy" when calculating gradients.

Although techniques such as clipping attempt to address off-policy discrepancies, our analysis reveals a deeper problem within GRPO: \textit{the objective itself is fundamentally ill-posed}. This deficiency becomes especially pronounced in scenarios involving large-scale models tasked with generating long-form responses, often resulting in severe collapse of the model. The root cause of the ill-posedness in the GRPO objective lies in the improper utilization of importance sampling weights. Importance sampling, in essence, seeks to compute the expectation of a function $f$ with respect to a target distribution $\pi_\text{tar}$ by appropriately re-weighting data sampled from a behavior distribution $\pi_\text{beh}$:

\begin{align}
\mathbb{E}_{z \sim \pi_\text{tar}} \left[ f(z) \right] = \mathbb{E}_{z \sim \pi_\text{beh}} \left[ \frac{ \pi_\text{tar} (z) }{ \pi_\text{beh} (z)} \, f(z) \right].
\end{align}

Importantly, this approach depends on aggregating a large number of samples ($N \gg 1$) drawn from the behavior distribution $\pi_\text{beh}$, allowing the importance weight $\frac{ \pi_\text{tar} (z) }{ \pi_\text{beh} (z)}$ to adequately adjust for discrepancies between the distributions.

Conversely, GRPO assigns the importance weight $\frac{ \pi_{\theta} (y_{i,t} | x, y_{i,<t}) }{ \pi_{\theta_\text{old}} (y_{i,t} | x,y_{i,<t})}$ independently at every token position $t$. This weighting relies solely on a single sample $y_{i,t}$ from each next-token distribution $\pi_{\theta_\text{old}}( \cdot | x, y_{i, <t})$, thereby failing to achieve proper correction between distributions as intended. Rather than mitigating bias, this procedure injects significant high-variance noise into the gradients used for training. Such variance accumulates, especially with longer sequences, and the issue is aggravated by the presence of clipping. Our experimental results indicate that these dynamics frequently result in catastrophic model collapse, a phenomenon that is typically irreversible. Once the model has collapsed, it is not possible to restore normal learning progress—even after resetting to earlier checkpoints, carefully readjusting hyperparameters such as clipping thresholds, increasing the sequence generation length, or altering RL querying strategies.

This observation highlights an essential flaw in the architecture of GRPO. The ineffectiveness of token-level importance weighting underscores a central tenet: \textbf{the granularity of the optimization target must align with that of the reward signal}. Because rewards are distributed across whole sequences, implementing off-policy correction at the token scale introduces complications. Consequently, we are driven to abandon token-based objectives in favor of leveraging importance weights and conducting optimization at the \textit{sequence level} instead.

\section{Algorithm}

\subsection{GSPO: Group Sequence Policy Optimization}

Although the token-level importance weight $\frac{ \pi_{\theta} (y_{i,t} | x, y_{i,<t}) }{ \pi_{\theta_\text{old}} (y_{i,t} | x,y_{i,<t})}$ presents challenges within GRPO, it is noteworthy that, in the realm of language generation, the \textit{sequence-level} importance weight $\frac{ \pi_{\theta} (y | x) }{ \pi_{\theta_\text{old}} (y | x)}$ possesses a well-defined theoretical significance. Specifically, it quantifies the extent to which the response $y$, drawn from $\pi_{\theta_\text{old}} (\cdot | x)$, differs from the distribution induced by $\pi_{\theta} (\cdot | x)$. This correspondence aligns directly with the sequence-level reward framework and provides an intuitive and informative signal for interpreting the clipping operation.

Drawing from this basic insight, we introduce the \textbf{Group Sequence Policy Optimization (GSPO)} algorithm. GSPO utilizes the following optimization objective at the sequence level:

\begin{align}
\mathcal{J}_\text{GSPO} (\theta) =
\mathbb{E}_{ x \sim \mathcal{D},\, \{y_i\}_{i=1}^G \sim \pi_{\theta_\text{old}}( \cdot | x) }
\left[ 
\frac{1}{G} \sum_{i=1}^{G}
\min \left( s_{i}(\theta)  \widehat{A}_{i},  \, \mathrm{clip} \left( s_{i}(\theta), 1 - {\varepsilon}, 1 + {\varepsilon} \right) \widehat{A}_{i} \right) 
\right],
\label{equ:gspo}
\end{align}

where we utilize the advantage estimation based on groups:

\begin{align}
\widehat{A}_{i} = \frac{r(x, y_i) - \mathrm{mean} \left( \{ r(x, y_i) \}_{i=1}^G \right) }{ \mathrm{std} \left( \{ r(x, y_i) \}_{i=1}^G \right) },
\end{align}

and construct the importance ratio $s_{i}(\theta)$ by utilizing the sequence likelihood as described in \citep{click}:

\begin{align}
s_{i}(\theta) = \left( \frac{ \pi_{\theta} (y_i | x) }{ \pi_{\theta_\text{old}} (y_i | x)} \right)^{\frac{1}{|y_i|}}
=
\exp \left( \frac{1}{|y_i|} \sum_{t=1}^{|y_i|} \log \frac{ \pi_{\theta} (y_{i,t} | x, y_{i,<t}) }{ \pi_{\theta_\text{old}} (y_{i,t} | x,y_{i,<t})} \right).
\end{align}

Consequently, GSPO performs clipping at the level of entire responses rather than on a per-token basis, thereby eliminating highly ``off-policy'' samples from gradient computation and ensuring consistency with sequence-level rewards and the optimization procedure. It is important to mention that length normalization is utilized in $s_{i}(\theta)$, which serves to decrease variance and to constrain $s_{i}(\theta)$ within a consistent numerical interval. Without this normalization, shifts in the likelihood of a handful of tokens could cause significant instability in the overall sequence-level importance ratio, necessitating the use of different clipping ranges for responses of varying lengths. Additionally, it is worth observing that the clipping ranges used in GSPO, as compared with prior methods (such as GRPO), are generally separated by several orders of magnitude, reflecting the unique definitions employed for importance ratios.

\subsection{Gradient Analysis}
\label{subsec:gradient}

The gradient of the GSPO objective can be obtained in the following way (for simplicity, we do not include clipping here):

\begin{align}
\nabla_{\theta} \mathcal{J}_\text{GSPO} (\theta)
=&\ 
\nabla_{\theta} \mathbb{E}_{ x \sim \mathcal{D},\, \{y_i\}_{i=1}^G \sim \pi_{\theta_\text{old}}( \cdot | x) }
\left[ 
\frac{1}{G} \sum_{i=1}^{G} s_{i}(\theta) \widehat{A}_{i}
\right] \\
= &\ 
\mathbb{E}_{ x \sim \mathcal{D},\, \{y_i\}_{i=1}^G \sim \pi_{\theta_\text{old}}( \cdot | x) }
\left[ 
\frac{1}{G} \sum_{i=1}^{G}
s_{i}(\theta) \widehat{A}_{i}
\cdot \nabla_{\theta} \log s_{i}(\theta)
\right] \\
=&\ 
\mathbb{E}_{ x \sim \mathcal{D},\, \{y_i\}_{i=1}^G \sim \pi_{\theta_\text{old}}( \cdot | x) }
\left[ 
\frac{1}{G} \sum_{i=1}^{G} 
\left( \frac{ \pi_{\theta} (y_i | x) }{ \pi_{\theta_\text{old}} (y_i | x)} \right)^{\frac{1}{|y_i|}} \widehat{A}_{i} 
\cdot \frac{1}{|y_i|} \sum_{t=1}^{|y_i|} \nabla_{\theta} \log \pi_{\theta} (y_{i,t} | x, y_{i,<t}) 
\right].
\label{equ:gspo_grad}
\end{align}

To facilitate comparison, the GRPO objective yields the following gradient (observe that $\widehat{A}_{i,t}$ is equivalent to $\widehat{A}_{i}$):

\begin{align}
\nabla_{\theta} \mathcal{J}_\text{GRPO} (\theta) 
=&\ 
\nabla_{\theta} \mathbb{E}_{ x \sim \mathcal{D},\, \{y_i\}_{i=1}^G \sim \pi_{\theta_\text{old}}( \cdot | x) }
\left[ \frac{1}{G} \sum_{i=1}^{G} \frac{1}{|y_i|} \sum_{t=1}^{|y_i|} 
w_{i,t}(\theta) \widehat{A}_{i,t}
\right] \\
=&\
\mathbb{E}_{ x \sim \mathcal{D},\, \{y_i\}_{i=1}^G \sim \pi_{\theta_\text{old}}( \cdot | x) }
\left[ \frac{1}{G} \sum_{i=1}^{G} \widehat{A}_{i} 
\cdot \frac{1}{|y_i|} \sum_{t=1}^{|y_i|} 
\frac{ \pi_{\theta} (y_{i,t} | x, y_{i,<t}) }{ \pi_{\theta_\text{old}} (y_{i,t} | x,y_{i,<t})} 
\nabla_{\theta} \log \pi_{\theta} (y_{i,t} | x, y_{i,<t})  
\right].
\end{align}

The principal difference between GSPO and GRPO pertains to \textit{the method by which each approach assigns weights to the gradients of token log likelihoods}. In GRPO, each token's contribution is scaled by its ``importance weight'' $\frac{ \pi_{\theta} (y_{i,t} | x, y_{i,<t}) }{ \pi_{\theta_\text{old}} (y_{i,t} | x,y_{i,<t})}$. These weights are not uniform and can range within $(0, 1+\varepsilon]$ when $\widehat{A}_{i} > 0$ or $[1-\varepsilon, +\infty)$ for $\widehat{A}_{i} < 0$. The non-trivial variation among these weights means their effects may accumulate during training, potentially resulting in instability or unpredictable training dynamics. Conversely, GSPO applies an equal weighting to all tokens within a given response, thereby removing the instability introduced by the unequal weighting in GRPO.

\subsection{GSPO-token: A Token-level Objective Variant}

In contexts such as multi-turn RL, a more detailed approach to advantage modulation than what is provided at the sequence level may be preferable. Therefore, we propose a token-level adaptation of GSPO, referred to as \textbf{GSPO-token}, which enables the assignment of advantages on a per-token basis:

\begin{align}
\mathcal{J}_\text{GSPO-token}(\theta)
=
\mathbb{E}_{ x \sim \mathcal{D},\, \{y_i\}_{i=1}^G \sim \pi_{\theta_\text{old}}( \cdot | x) }
\left[ 
\frac{1}{G} \sum_{i=1}^{G} \frac{1}{|y_i|} \sum_{t=1}^{|y_i|} 
\min \left( s_{i,t}(\theta) \widehat{A}_{i,t},  \, \mathrm{clip} \left( s_{i,t}(\theta), 1 - {\varepsilon}, 1 + {\varepsilon} \right) \widehat{A}_{i,t} \right) 
\right],
\label{equ:gspo-token}
\end{align}

in which

\begin{align}
s_{i,t}(\theta) 
= \mathrm{sg} \left[ s_{i}(\theta) \right]  \cdot \frac{ \pi_{\theta} (y_{i,t} | x, y_{i,<t}) }{ \mathrm{sg} \left[ \pi_{\theta} (y_{i,t} | x, y_{i,<t}) \right] },
\end{align}

Here, $\mathrm{sg}[\cdot]$ refers to extracting the numerical value while halting any gradient propagation, analogous to the \texttt{detach} function in PyTorch. The GSPO-token gradient can be computed as follows:

\begin{align}
\nabla_{\theta} \mathcal{J}_\text{GSPO-token}(\theta)
=&\ 
\nabla_{\theta} \mathbb{E}_{ x \sim \mathcal{D},\, \{y_i\}_{i=1}^G \sim \pi_{\theta_\text{old}}( \cdot | x) }
\left[ 
\frac{1}{G} \sum_{i=1}^{G}
\frac{1}{|y_i|} \sum_{t=1}^{|y_i|} 
s_{i,t}(\theta) \widehat{A}_{i,t}
\right] \\
=&\ 
\mathbb{E}_{ x \sim \mathcal{D},\, \{y_i\}_{i=1}^G \sim \pi_{\theta_\text{old}}( \cdot | x) }
\left[ 
\frac{1}{G} \sum_{i=1}^{G} s_i (\theta) \cdot \frac{1}{|y_i|} \sum_{t=1}^{|y_i|} 
\widehat{A}_{i,t} \frac{ \nabla_{\theta} \pi_{\theta} (y_{i,t} | x, y_{i,<t}) }{ \pi_{\theta} (y_{i,t} | x, y_{i,<t}) }
\right] \\
=&\ 
\mathbb{E}_{ x \sim \mathcal{D},\, \{y_i\}_{i=1}^G \sim \pi_{\theta_\text{old}}( \cdot | x) }
\left[ 
\frac{1}{G} \sum_{i=1}^{G}  \left( \frac{ \pi_{\theta} (y_i | x) }{ \pi_{\theta_\text{old}} (y_i | x)} \right)^{\frac{1}{|y_i|}}
\cdot \frac{1}{|y_i|} \sum_{t=1}^{|y_i|}
\widehat{A}_{i,t} \nabla_{\theta} \log \pi_{\theta} (y_{i,t} | x, y_{i,<t})  
\right].
\label{equ:gspo-token_grad}
\end{align}

It should be observed that $\frac{ \pi_{\theta} (y_{i,t} | x, y_{i,<t}) }{ \mathrm{sg} \left[ \pi_{\theta} (y_{i,t} | x, y_{i,<t}) \right] }$ evaluates to 1, ensuring that $s_{i,t}(\theta)$ is, in fact, equal in value to $s_{i}(\theta)$ from a numerical perspective. By examining Equation~\eqref{equ:gspo} alongside~\eqref{equ:gspo-token}, as well as Equation~\eqref{equ:gspo_grad} with~\eqref{equ:gspo-token_grad}, it is evident that GSPO-token yields the same numerical results as GSPO in terms of the objective function for optimization, the condition for clipping, and the formal expression for the gradient, provided the advantage values for all tokens in a sequence $y_i$ are held constant (that is, $\widehat{A}_{i,t} = \widehat{A}_{i}$). Nonetheless, GSPO-token introduces the added capability of independently tuning the advantage for each individual token.

\section{Experiments and Discussion}

\subsection{Empirical Results}

We conduct experiments utilizing a cold-start model that has been fine-tuned from Qwen3-30B-A3B-Base. The results include plots of both training reward progression and evaluation performance on several benchmarks: AIME'24 (average Pass@1 across 32 samples), LiveCodeBench (202410-202502, average Pass@1 across 8 samples), and CodeForces (Elo Rating). Throughout the reinforcement learning training process, each rollout data batch is subdivided into four mini-batches to facilitate gradient updates. For GSPO, we configure the left and right clipping bounds in Equation~\eqref{equ:gspo} to be 3e-4 and 4e-4, respectively. For comparative purposes, GRPO is adopted as a baseline; its left and right clipping thresholds in Equation~\eqref{equ:grpo} are set to 0.2 and 0.27, chosen through careful tuning for an equitable assessment. It should be emphasized that GRPO requires the use of the Routing Replay strategy to guarantee reliable MoE RL convergence, a topic we further elaborate on in \S~\ref{subsec:routing_replay}, whereas \textit{GSPO eliminates this requirement}.

As illustrated in Figure~\ref{fig:results}, training with GSPO exhibits stable progress over time. Notably, \textbf{GSPO consistently enhances performance as additional training compute is allocated, the query set is periodically refreshed, and the generation length is increased}. In addition, GSPO outperforms GRPO in terms of training efficiency, achieving higher accuracy and superior benchmark results while utilizing equivalent training compute and query consumption. Ultimately, we have implemented GSPO in the reinforcement learning training of the latest Qwen3 models, which compellingly demonstrates GSPO’s effectiveness in harnessing the scaling capabilities of RL for large language models.

\subsection{Curious Observation on Clipping Fractions}

An important difference between GSPO and GRPO lies in the approach to clipping: GSPO applies clipping at the response level, as opposed to the token-level clipping employed by GRPO. Notably, as illustrated in Figure~\ref{fig:clipping}, there is a discrepancy of approximately two orders of magnitude in the proportion of clipped tokens between the two methods, and modifying the clipping thresholds does not diminish this significant gap. Despite the fact that GSPO discards a considerably larger share of tokens—resulting in fewer tokens being available for training or gradient calculation—it consistently outperforms GRPO in terms of training efficiency. This somewhat surprising outcome—that increased clipping correlates with greater training efficiency—suggests that GRPO's method of estimating gradients at the token level introduces substantial noise and is less effective for making use of sampled data. Conversely, by leveraging sequence-level information, GSPO yields a learning signal that is both more stable and more beneficial.

\subsection{Benefit of GSPO for MoE Training}
\label{subsec:routing_replay}

\paragraph{Background}
While RL training of dense architectures generally proceeds in a stable manner, the inherent sparsity of MoE models brings forth distinct challenges. Specifically, our observations indicate that applying the GRPO algorithm to MoE models results in pronounced \textit{expert-activation volatility}, which disrupts the convergence of RL training. More concretely, following one or more rounds of gradient descent, the selection of experts engaged for generating an identical output often changes substantially. For instance, when experimenting with the 48-layer Qwen3-30B-A3B-Base model, we found that each RL gradient update leads to approximately 10\% of the experts for a given rollout sample under the new policy $\pi_{\theta}$ differing from those chosen under the previous policy $\pi_{\theta_\text{old}}$. This effect intensifies in MoE models with greater depth, causing the token-level importance weights $w_{i,t}(\theta) = \frac{ \pi_{\theta} (y_{i,t} | x, y_{i,<t}) }{ \pi_{\theta_\text{old}} (y_{i,t} | x,y_{i,<t})}$ to vary erratically, thereby undermining their reliability—as further elaborated in \S~\ref{sec:motivation} and~\ref{subsec:gradient}—and ultimately impeding the typical progression of RL convergence.

\paragraph{Our Previous Approach} In our earlier work, we addressed this issue by introducing the \textbf{Routing Replay} training method. This approach involves saving the expert routing choices determined by $\pi_{\theta_\text{old}}$ and subsequently applying these same routing selections in $\pi_{\theta}$ when evaluating the importance ratios $w_{i,t}(\theta) = \frac{ \pi_{\theta} (y_{i,t} | x, y_{i,<t}) }{ \pi_{\theta_\text{old}} (y_{i,t} | x,y_{i,<t})}$. By synchronizing the network activations for every token $y_{i,t}$ between $\pi_{\theta}$ and $\pi_{\theta_\text{old}}$, we enhance the reliability of token-wise importance ratios, thus maintaining consistent optimization of the shared network through each gradient step. As illustrated in Figure~\ref{fig:routing_replay}, Routing Replay proves critical for achieving stable convergence during GRPO training in MoE models.

\paragraph{Benefit of GSPO} While Routing Replay facilitates the successful GRPO training of MoE models, its approach—relying on the repeated use of prior routing choices—incurs significant memory and communication costs, and further restricts the attainable model capacity. In comparison, as depicted in Figure~\ref{fig:results}, GSPO avoids dependence on Routing Replay altogether. It can directly calculate the importance ratios $s_i(\theta)$ as standard, leading to steady convergence and optimization during training. The main conceptual advance of GSPO lies in its exclusive consideration of sequence-level likelihoods, $\pi_{\theta} (y_i | x)$, as opposed to token-level probabilities, $\pi_{\theta} (y_{i,t} | x, y_{i,<t})$, thus rendering it robust to token-level fluctuations. Given that the language modeling strength of the MoE framework remains intact, the overall sequence likelihood stays comparatively stable. To conclude, GSPO directly tackles and mitigates the problem of expert-activation instability inherent to MoE models, thus rendering workarounds like Routing Replay unnecessary. As a result, GSPO streamlines and fortifies the training process and enables exploitation of the full expressive power of the MoE architecture without imposed limits.

\subsection{Benefit of GSPO for RL Infrastructure}

Due to the differences in numerical precision between training platforms such as Megatron and inference frameworks like SGLang and vLLM, it is common practice to recalculate the likelihoods of sampled responses under the prior policy $\pi_{\theta_\text{old}}$ using the training engine. Nevertheless, GSPO relies on sequence-level rather than token-level likelihoods during optimization. Conceptually, sequence-level likelihoods are significantly more robust to variations in numerical precision. As a result, GSPO allows the likelihoods produced by the inference engine to be used directly in the optimization process, eliminating the necessity for recomputation via the training engine. This feature is particularly advantageous for tasks involving partial rollout, multi-turn reinforcement learning, or in frameworks where training and inference are decoupled.

\section{Conclusion}

We introduce Group Sequence Policy Optimization (GSPO), a novel reinforcement learning technique designed to train large language models. Adhering to the principles of importance sampling, GSPO constructs importance ratios via sequence likelihood and applies operations such as sequence-level clipping, reward assignment, and optimization. Empirically, GSPO surpasses GRPO in terms of training stability, efficiency, and overall performance, and proves particularly effective for large-scale RL applications in MoE model training—contributing substantially to the significant advancements observed in recent Qwen3 models. As a robust and scalable framework, GSPO will serve as the algorithmic basis for ongoing reinforcement learning scaling, which we anticipate will drive major progress in the development of advanced intelligence.

\end{document}